% ==================== THEME 1 — VARIABLES SIMPLES (entraînement) ====================

% ---- var-declaration ----
\begin{question}{ent-var-declaration}
    Quelle ligne déclare correctement deux variables entières \texttt{a} et \texttt{b} en une seule instruction ?
    \begin{multicols}{4}
        \begin{reponses}
            \bonne{\UPSTIprofOnly[1]{\checkmark\ }\lstinline[language=C]|int a, b;|}
            \mauvaise{\lstinline[language=C]|int a; int b|}
            \mauvaise{\lstinline[language=C]|int a - b;|}
            \mauvaise{\lstinline[language=C]|int (a, b);|}
        \end{reponses}
    \end{multicols}
\end{question}

% ---- var-type-decimal ----
\begin{question}{ent-var-type-decimal}
    Un capteur de température renvoie une mesure telle que \texttt{21.75}. Quel type utiliser en priorité pour la stocker en C ?
    \begin{multicols}{4}
        \begin{reponses}
            \mauvaise{\lstinline[language=C]|int|}
            \bonne{\UPSTIprofOnly[1]{\checkmark\ }\lstinline[language=C]|float|}
            \mauvaise{\lstinline[language=C]|char|}
            \mauvaise{\lstinline[language=C]|string|}
        \end{reponses}
    \end{multicols}
\end{question}

% ---- var-type-caractere ----
\begin{question}{ent-var-type-caractere}
    On veut stocker la réponse donnée à un QCM, une seule lettre parmi \texttt{'A'}, \texttt{'B'}, \texttt{'C'}, \texttt{'D'}. Quel type utiliser ?
    \begin{multicols}{4}
        \begin{reponses}
            \mauvaise{\lstinline[language=C]|int|}
            \mauvaise{\lstinline[language=C]|float|}
            \bonne{\UPSTIprofOnly[1]{\checkmark\ }\lstinline[language=C]|char|}
            \mauvaise{\lstinline[language=C]|void|}
        \end{reponses}
    \end{multicols}
\end{question}

% ---- var-nommage ----
\begin{question}{ent-var-nommage}
    Quel nom de variable est \textbf{invalide} en C ?
    \begin{multicols}{5}
        \begin{reponses}
            \mauvaise{\lstinline[language=C]|nb_points|}
            \mauvaise{\lstinline[language=C]|_etat|}
            \bonne{\UPSTIprofOnly[1]{\checkmark\ }\lstinline[language=C]|return|}
            \mauvaise{\lstinline[language=C]|joueur1|}
            \mauvaise{\lstinline[language=C]|ScoreMax|}
        \end{reponses}
    \end{multicols}
\end{question}

% ---- var-initialisation ----
\begin{question}{ent-var-initialisation}
    Une variable entière non initialisée (\lstinline[language=C]|int n;|) contient-elle nécessairement la valeur \texttt{0} ?
    \begin{multicols}{2}
        \begin{reponses}
            \bonne{\UPSTIprofOnly[1]{\checkmark\ }Non, sa valeur est indéterminée (le contenu dépend de ce qu'il y avait avant en mémoire)}
            \mauvaise{Oui, toute variable C vaut \texttt{0} par défaut}
            \mauvaise{Non, elle contient toujours \texttt{-1}}
            \mauvaise{Non, le programme refuse de compiler}
        \end{reponses}
    \end{multicols}
\end{question}
